% Authority: EGA I chapter1-frontmatter-fr.tex, SHA-256 7B2D0F8F812EBA3121202F0AE6415FFC6C281B8428DA8F0F72D89DF1CEC01708.
\clearpage
\oldpage[I]{79}
\phantomsection
\label{chapter:I-ko}

\begin{center}
{\large 제1장\par}
\vspace{1.5em}
{\Large\bfseries 스킴의 언어\par}
\vspace{1.8em}
{\bfseries 목차\par}
\end{center}

\medskip
\begin{tabular}{@{}r@{\ }l@{}}
\S~1.  & 아핀 스킴.\\
\S~2.  & 준스킴과 준스킴의 사상.\\
\S~3.  & 준스킴의 곱.\\
\S~4.  & 부분준스킴과 몰입 사상.\\
\S~5.  & 축소 준스킴; 분리 조건.\\
\S~6.  & 유한성 조건.\\
\S~7.  & 유리 사상.\\
\S~8.  & 슈발레 스킴.\\
\S~9.  & 준연접층에 관한 보충.\\
\S~10. & 형식 스킴.
\end{tabular}

\medskip
\S\S~1에서~8까지는 이후 전체에서 사용할 언어를 전개하는 것에서
크게 벗어나지 않는다. 다만 이 논고의 전반적인 정신에 따라
\S\S~7과~8은 다른 절들보다 덜 쓰이고, 쓰이더라도 덜 본질적이라는
점을 밝혀 둔다. 슈발레 스킴을 다룬 것도 슈발레~\cite{I-1}와
나가타~\cite{I-9}의 언어와 연결하기 위해서일 뿐이다. \S~9에서는
준연접층에 관한 정의와 결과들을 제시한다. 그 가운데에는 알려진
가환대수의 개념들을 ``기하학적'' 언어로 옮기는 데 그치지 않고 이미
전역적인 성격을 띠는 것도 있으며, 이러한 결과들은 바로 다음 장들부터
사상을 전역적으로 연구하는 데 반드시 필요하다. 끝으로 \S~10에서는
스킴 개념의 한 일반화를 도입한다. 이는 제III장에서 고유 사상의 코호몰로지
연구에 관한 기본 결과들을 편리하게 서술하고 증명하는 중간 도구로 쓰일
것이다. 한편 형식 스킴의 개념은 ``모듈라이 이론''의 몇몇 사실, 곧
대수다양체의 분류 문제를 표현하는 데 없어서는 안 될 것으로 보인다.
\S~10의 결과들은 제III장 \S~3 이전에는 사용하지 않으므로, 그때까지는
이 절을 건너뛰어 읽기를 권한다.

\clearpage
